\chapter{Dressing: Placing Rather Than Sprinkling}
\label{ch:dressing}

\section{The Central Decision, and the Rewrite That Produced It}

Dressing is the third pass over a realized world, after the structures are stamped
and the terrain is painted. Where the painter adds no geometry and only rewrites
materials, this pass exists to add geometry: trees, boulders, flora, paths,
buildings.

The obvious design for such a pass is a set of density fields assigned to shapes,
the way a terrain theme is assigned. That was the first implementation, and it was
replaced. The reasoning is the clearest statement in the system of why map-making
resists automation:

\begin{quote}
A tree is cover and a boulder is a wall, so where each one stands decides how the
map plays, and a decision about how a map plays belongs to the person making it
rather than to a noise field. So every part of this stage is a thing put somewhere:
a route is dragged, an area of cover is traced, a tree and a rock are clicked into
place, and each carries its own knobs. The fields that remain are the ones
\emph{inside} a drawn area, where placing them one at a time would be data entry
rather than authoring.
\end{quote}

That line, between a decision that affects play and a decision that is merely
tedious, is drawn once and then holds. Which tree stands at the mouth of a lane is
authored. Which blade of grass fills a square metre of meadow is not.

\section{Running Last Is What Makes It Tractable}

The pass reads the finished, painted world. That ordering is load-bearing rather
than incidental: the painter has already decided, per cell, what the surface
\emph{is}, so soil accepts flora and quartz does not, grass can be replaced by a
stroke and a monument's wool cannot. A pass running earlier would have to predict
all of that.

The one elevation model it needs is the first air height above each column, and the
one it is given is the tops of everything on the board \emph{that is not a made
thing}.

\begin{ruling}{A Made Thing Is Not Ground}
Nothing that rests on the board treats a made thing as ground. The failure that
established this is instructive: a cloud drawn at $y=78$ over a car park was the
top of every column beneath it, so the goal there read the cloud and was stamped at
$y=83$, over a build ceiling of 68. The same elevation feeds the structure-site
check, which would otherwise report a shed under a balloon as standing on a
fifty-block plinth.
\end{ruling}

Two consequences follow for a stacked board. A prop carries an optional layer and
rests on that layer's own surfaces; naming a layer the board has no ground on is
declined. And the claim book that stops two props sharing a cell is keyed on the
layer as well as the cell, so two props share a cell only where they share a layer.

\section{Keeping Ground Clear}

A shape can declare that it is not ground to dress. A shape marked \id{keepClear}
puts its own columns into the mask exactly and with no margin, so a road still runs
to a gate while the wall either side of it keeps its top course.

The failure without it is worth stating because it is not obvious. A stroke
repaints whatever it crosses. A channel, whose water line is the \emph{lowest}
surface its band crosses, cuts every other column in the band down to that line.
On a wall standing seventeen courses over a river, that is a hole through the wall
rather than a bank.

The keep-out mask itself gathers spawns, wool rooms, stated structures, built
columns, the shapes that marked themselves clear, and every door's approach. It is
not the map contract's protection regions, and destroyables and cores are
deliberately not in it.

\section{Six Primitives}

\begin{description}[leftmargin=1.6em,style=nextline,itemsep=5pt]
  \item[Noise] The same deterministic hash-from-cell the terrain patterns use, never
    a random number generator. A density field decides where flora goes, a
    low-frequency field gathers flowers into meadows and trees into groves, and a
    per-cell value is the dice a worn stroke rolls.
  \item[Mask] Eligibility, read from the painted surface and from the keep-out set.
    The whole of it applies to anything that stands on ground; only the structure
    part applies to a stroke, which stands on nothing.
  \item[Placement] A \emph{point} for props that stand somewhere, a \emph{drawn
    outline} for ones that cover a stretch, and a \emph{dragged rectangle} for the
    one whose stamp takes a footprint. A marker seats on the ground and can
    therefore only be dropped on rasterized terrain; an area has no such limit.
  \item[Reshaping] A selected prop wears one grip per point. Reshaping never changes
    how many points a prop has.
  \item[Stamp] A three-dimensional volume seated on the surface and written cell by
    cell. A prop seats on the \emph{lowest} column of its own footprint, so it sits
    into a slope rather than floating over the low side.
  \item[The fan] Every prop is placed once and stamped at every image of its orbit,
    in the prop's own local frame, with each offset turned by that image's
    transform.
\end{description}

Two details in that list repay attention.

\begin{ruling}{Only the Feet Are Asked About}
What a prop \emph{rests} on needs real ground; what it merely reaches over does
not. A cell of the prop standing where something already is is skipped, and nothing
else happens.
\end{ruling}

The fan carries a mistake worth naming, because it is invisible in a top-down
render and obvious in play. Mirroring only the anchor leaves both teams with the
same unmirrored prop shape, so a boulder with a lobe to its east has a lobe to its
east on both sides of the map. The stamp is also all-or-nothing per image: if one
image's ground is missing or kept clear, that image is skipped rather than clipped.

\section{Clipping, and the Piece of Tree Floating Next to a Building}

A prop stamped against an obstruction can end in one of three states, and the
studio distinguishes them because only one is a fault.

\textbf{Truncated} is a prop with a face taken off it. Measured against a wall two
blocks off its centre, an erratic of size seven loses 36\% of its 1{,}089 blocks
and severs none of them. That is acceptable: it reads as a boulder against a wall.

\textbf{Cut in two} is a prop whose obstruction took out the blocks joining a piece
of it to the rest. Nothing declines, the tree is in the world, and there is a piece
of tree floating next to a building.

The measurement is taken against the prop \emph{as it would have stood on open
ground}, and connectivity is judged by face adjacency, because a block joined to
its neighbour at a corner alone has air on all six of its own faces and is seen
straight past.

\section{Placement Is a Query}

\Cref{ch:division} and \cref{ch:session} both arrive at this point from different
directions, and it belongs here as mechanism. The studio can answer where a prop or
a building may legally stand, and an agent that asks converges immediately where an
agent that guesses does not.

Across three separate runs, declined props on successive dressing passes ran seven,
four, two, zero; then six, five, three, none; then five, five, five, six. In every
case the pass that reached zero was the one where every position came off the
studio's own seats mask rather than off the agent's judgement. One run states the
final tally plainly: 78 placed, 0 declined.

The same read answers a question about the board rather than about the prop. Asked
whether a composed board had room for buildings, the seats read answered that a
12$\times$8 house seated in \emph{zero} cells, a 9$\times$7 in zero, and a
7$\times$5 in seventy. That is a fact about the board's grain, and no picture of it
would have said so.

\section{Trees That Are Copied Rather Than Grown}

A tree may be a recipe the studio grows, or it may be \emph{copied}: a body
recorded block for block out of a real world and filed in the library under a name.
A placement then names that style. The distinction matters for the look of a board
more than almost anything else in this chapter, because a grown tree is
recognisably generated and a copied one is recognisably built.

One caution from the corpus, which no gate reports: nothing asks what block a prop
is seated \emph{on}. The site check asks whether there is ground under a prop, not
what that ground is. One run found thirty-five of ninety-eight trees not seated on
soil, including trees on ice, on coal ore, six on a ruin's masonry, and one planted
in another tree's leaves. The only read that answers it is one column per prop.
